%% Separately authored finite-CNF certificate construction; no inherited text replaced.
\chapter{DRAT Certificates and Checked Unsatisfiability}
\label{sec:p12-drat-certificates-and-checking}

The certificate boundary of
Definition~\ref{def:p12-sat-certificate-boundary} becomes a mathematical
guarantee only after the accepted log operations have been specified and proved
sound.  This chapter supplies that guarantee for finite propositional CNF.
The source presentation of Wetzler, Heule, and Hunt is the 2014 DRAT-trim
paper: its asymmetric-literal-addition definition, its pivot-retaining RAT
test, its clause-set update, and its first-literal pivot convention are retained
below \cite[Secs.~2--3]{WetzlerHeuleHunt2014AuthorDRAT}.  The separately dated
2016 format description has a multiset interpretation and additional parser
and deletion conventions.  Those differences are stated in
Section~\ref{sec:p12-drat-format-versions}, not attributed retrospectively to
the 2014 paper.  The proofs in the annex are complete mathematical
reconstructions of the checking core, not a verification of an executable
DRAT-trim implementation.

\section{Clauses, assumptions, and unit propagation}

\begin{hbdefinition}{def:p12-drat-finite-cnf}{Finite-CNF semantic interface}
Fix a finite variable set \(X\).  A literal \(l\) is \(x\) or \(\neg x\)
for \(x\in X\); its complement \(\bar l\) interchanges those two forms, so
\(\overline{\bar l}=l\).  A clause \(C\) is a finite set of literals,
interpreted as their disjunction.  In particular \(\varnothing\), when typed
as a clause, is false.  A clause containing both \(l\) and \(\bar l\) is
\emph{tautological} and is true under every valuation.  A CNF database \(F\)
is a finite set of clauses, interpreted conjunctively; the empty database is
true.  A valuation \(v:X\to\{0,1\}\) satisfies a clause when one of its
literals is true, and satisfies \(F\), written \(v\models F\), when it
satisfies every clause of \(F\).  The proposition \(\operatorname{Sat}(F)\)
means that such a valuation exists.  These types distinguish an empty clause
from an empty database and literal complement from negation of a CNF.  They
are the input types used in every checking result below; no restriction to
Horn clauses, bounded width, or a fixed clause count is imposed.

Repeated occurrences of a literal have no semantic effect and are represented
once by this mathematical clause type.  A serialized clause remains an
ordered list, however: its first literal specifies the RAT pivot.  Retaining
that first literal is necessary even when the remaining list is normalized.
The mathematical interface admits tautological clauses in order to state its
boundary cases; it is not permission to submit a tautological clause to a
format version whose grammar restrictions forbid it.
\index{CNF!certificate semantics}\index{clause!empty}
\end{hbdefinition}

\begin{hbdefinition}{def:p12-drat-up}{Trails and the conflict test}
A \emph{trail} \(\tau=(l_1,\ldots,l_k)\) is a sequence of signed
assignments to distinct variables.  It determines the partial valuation that
makes every listed literal true.  A finite assumption set \(A\) initializes
the trail; if it contains complementary literals, the check immediately
returns the conflict token \(\bot\).  Otherwise list its literals in any
order.  For a database clause \(D\), first ignore it if a trail literal makes
it true.  If it is not already true, its residual
\(\operatorname{res}_{\tau}(D)\) consists of the literals not made false
by the trail.  An empty residual is a conflict.  A singleton residual
\(\{l\}\) permits appending the previously unassigned literal \(l\).
Repeat until a conflict occurs or no singleton residual remains.  Write
\(\operatorname{UP}(F,A)=\bot\) for the conflict outcome.  A saturated
nonconflicting trail is only a failed conflict test, not a SAT certificate.

The empty input clause conflicts before any propagation; the empty database
does not conflict under consistent assumptions.  A tautological clause can
neither conflict nor force a literal: it is either already satisfied or still
has both unassigned polarities.  These conventions prevent the disappearance
of an already satisfied clause from being confused with a falsified clause.
The trail here has assumptions and forced assignments, but no solver decision
levels, restarts, or backjump operation.  It is needed solely to make each
local redundancy test explicit.
\index{trail!certificate checking}\index{unit propagation!certificate checking}
\end{hbdefinition}

\begin{hblemma}{thm:p12-drat-up-sound}{proved=annex:p12-drat-up-sound}{Unit-propagation termination and conflict soundness}
For every finite \(F,X,A\) as above, the conflict test terminates.  Every
model of \(F\) satisfying the initial assumptions satisfies every literal
subsequently propagated.  If \(\operatorname{UP}(F,A)=\bot\), no model
of \(F\) satisfies all assumptions in \(A\).
\end{hblemma}

\section{The 2014 asymmetric-tautology presentation}

\begin{hbdefinition}{def:p12-drat-ala-rup}{ALA, AT, and reverse unit propagation}
For a clause \(C\), the 2014 operation
\(\operatorname{ALA}(F,C)\) is the least clause containing \(C\) closed
under asymmetric literal addition: whenever
\(\{l_1,\ldots,l_k,l\}\in F\) and all
\(l_1,\ldots,l_k\) belong to the current clause, add \(\bar l\).
Here \(k\) may be zero.  Equivalently, for each \(D\in F\) and
\(l\in D\), add \(\bar l\) when
\(D\setminus\{l\}\) is contained in the current clause.  Stop at the
finite least fixed point; the order of eligible additions does not change it.
The source calls \(C\) an \emph{asymmetric tautology} (AT) when that
closure contains complementary literals
\cite[Def.~1]{WetzlerHeuleHunt2014AuthorDRAT}.

The operational predicate \(\operatorname{RUP}(F,C)\), called
\emph{reverse unit propagation}, means
\[
 \operatorname{UP}(F,\bar C)=\bot,
 \qquad \bar C=\{\bar l:l\in C\}.
\]
Thus \(\bar C\) is a set of simultaneously true assumption literals,
equivalently a conjunction of unit assumptions, not a clause asserting their
disjunction.  When \(C=\varnothing\), it is the empty assumption set and
the test is ordinary unit refutation of \(F\).  A tautological \(C\)
passes immediately because its assumptions are inconsistent.

For databases without an empty clause, the source's AT and this RUP predicate
coincide.  If \(F\) already contains the empty clause, conflict must instead
be recognized directly: the literal-extension rule has no literal to select
from that clause.  This explicit terminal case completes the checking
interface at arbitrary finite CNF; it is not an unstated alteration of the
source's literal-extension formula.  These names are needed to distinguish a
finite syntactic conflict check from semantic entailment, and are used in the
RAT tests and the trace induction.
\index{asymmetric literal addition}\index{asymmetric tautology}
\index{reverse unit propagation}
\end{hbdefinition}

\begin{hblemma}{thm:p12-drat-ala-rup}{proved=annex:p12-drat-ala-rup}{ALA--RUP correspondence and monotonicity}
The ALA operation is well defined.  If \(F\) has no empty clause, then
\(\operatorname{ALA}(F,C)\) is tautological if and only if
\(\operatorname{RUP}(F,C)\).  Every RUP clause is entailed by \(F\).
Moreover, if \(F\subseteq G\), \(C\subseteq E\), and
\(\operatorname{RUP}(F,C)\), then \(\operatorname{RUP}(G,E)\).
The entailment and monotonicity statements include empty clauses and
tautological clauses under the operational definition.
\end{hblemma}

\section{RAT, resolvents, and model repair}

\begin{hbdefinition}{def:p12-drat-rat}{The exact pivot-retaining RAT test}
Let \(C\ne\varnothing\) and choose \(l\in C\).  For every current
clause \(D\in F\) containing \(\bar l\), the ordinary resolvent and the
2014 test clause are respectively
\[
 \operatorname{Res}_{l}(C,D)
   =(C\setminus\{l\})\cup(D\setminus\{\bar l\}),
 \qquad
 C\cup(D\setminus\{\bar l\}).
\]
The \emph{resolution asymmetric tautology} predicate
\(\operatorname{RAT}(F,C,l)\) requires
\[
 \operatorname{RUP}\bigl(F,C\cup(D\setminus\{\bar l\})\bigr)
 \quad\text{for every }D\in F\text{ with }\bar l\in D.
\]
This is the pivot-retaining formula of the 2014 paper and also of the 2016
format note \cite[Sec.~2]{WetzlerHeuleHunt2014AuthorDRAT}
\cite[Clause Addition Preconditions]{Heule2016DRATFormatV1}.
The clause \(C\) has RAT with respect to \(F\) exactly when some
\(l\in C\) satisfies this predicate.  The quantifier ranges over the entire
active database, not merely the clauses
chosen by a solver or supplied as a partial hint list.  If there are no
opposing clauses, the condition is vacuous.  An empty added clause has no pivot
and cannot use this vacuity; it must pass RUP.  A tautological test clause
passes by inconsistent assumptions; a non-tautological test clause needs an
actual unit conflict.  RAT is a clause-redundancy condition preserving the
existence of a model, not a claim that the added clause is already true in
every old model.

The ordinary-resolvent formulation can be used only with its side conditions
\(l\in C\), \(\bar l\in D\), and \(D\in F\) in force.  The following
bridge justifies it; the name ``resolvent'' alone does not justify dropping the
pivot from the source test.
\index{resolution asymmetric tautology}\index{RAT!pivot}
\end{hbdefinition}

\begin{hblemma}{thm:p12-drat-resolvent-test}{proved=annex:p12-drat-resolvent-test}{Pivot-retaining and ordinary-resolvent tests}
If \(l\in C\), \(\bar l\in D\), and \(D\in F\), then
\[
 \operatorname{RUP}\bigl(F,C\cup(D\setminus\{\bar l\})\bigr)
 \quad\text{if and only if}\quad
 \operatorname{RUP}\bigl(F,\operatorname{Res}_{l}(C,D)\bigr).
\]
In particular every nonempty RUP clause is RAT on each of its literals.
Neither statement gives an empty clause a pivot.
\end{hblemma}

\begin{hbtheorem}{thm:p12-drat-rat-preservation}{proved=annex:p12-drat-rat-preservation}{RAT preserves satisfiability by pivot repair}
For every finite database \(F\), nonempty clause \(C\), and pivot
\(l\in C\), if \(\operatorname{RAT}(F,C,l)\), then
\[
 \operatorname{Sat}(F)\quad\text{if and only if}\quad
 \operatorname{Sat}(F\cup\{C\}).
\]
More precisely, each model \(v\) of \(F\) already satisfying \(C\) is
retained; a model falsifying \(C\) can be repaired by changing only the
pivot variable so that \(l\) becomes true.  The repaired valuation satisfies
every old clause and \(C\).  This statement permits fresh variables in
\(C\) once valuations are extended to the finite combined variable set.
It is an equisatisfiability theorem, not formula equivalence or preservation
of the set or number of models.
\end{hbtheorem}

\begin{hbexample}{def:p12-drat-nonrup-rat}{A source RAT addition that is not RUP}
The 2014 source example has
\[
 F=\{\{\neg a,b\},\{a,c\},\{b,\neg c\}\},\qquad C=\{a\}.
\]
To test RUP, assume \(\neg a\).  The clause \(\{a,c\}\) forces
\(c\), and \(\{b,\neg c\}\) forces \(b\).  All clauses are then
satisfied; there is no conflict.  Indeed \(a=0,b=1,c=1\) is a model of
\(F\) falsifying \(C\), so entailment itself fails.

For RAT on \(a\), the only opposing clause is \(\{\neg a,b\}\).
Its source test clause is \(\{a,b\}\), retaining the pivot.  Assume
\(\neg a,\neg b\): \(\{a,c\}\) forces \(c\), whereas
\(\{b,\neg c\}\) forces \(\neg c\), producing a conflict.  Therefore
\(C\) is RAT and can be added.  Repairing the displayed countermodel
changes only \(a\) to \(1\); the resulting model still satisfies all three
old clauses.  The source also identifies \(\{b\}\) as AT: assuming
\(\neg b\) forces \(\neg a\) from the first clause and \(\neg c\)
from the third, falsifying \(\{a,c\}\)
\cite[Example~1]{WetzlerHeuleHunt2014AuthorDRAT}.
This example witnesses the distinction needed in the general preservation
theorem; it is not substituted for that theorem's proof.
\end{hbexample}

\section{Trace construction and the independent checking guarantee}

\begin{hbdefinition}{def:p12-drat-trace}{The 2014 set-based addition/deletion trace}
A proof log \(P\) is a finite ordered sequence of clause additions and
clause deletions.  In the 2014 clause-set presentation, with \(L_i\) the
clause at step \(i\), define the active database by
\[
 F^0_P=F,\qquad
 F^i_P=
 \begin{cases}
  F^{i-1}_P\setminus\{L_i\},&\text{a deletion step},\\
  F^{i-1}_P\cup\{L_i\},&\text{an addition step}.
 \end{cases}
\]
The source expects the deleted clause to be present.  A strict checker may
reject an absent deletion; permitting it as a no-op is a separately declared
policy and cannot manufacture a conflict.  No redundancy test is required
for deletion.  Its guarantee is only
\(\operatorname{Sat}(F^{i-1}_P)\Rightarrow\operatorname{Sat}(F^i_P)\).
For example deleting \(\{\neg p\}\) from
\(\{\{p\},\{\neg p\}\}\) produces a satisfiable database.
Thus deletion is not an equisatisfiability step in general.

Before inserting an addition, test RUP against \(F^{i-1}_P\).  If that
fails and the clause is nonempty, test RAT on the first literal of its
serialized list against that same database.  Reject if both tests fail.
The final empty-clause addition must pass RUP against the database
\emph{before} it is inserted.  Testing after insertion would make every such
step trivially pass.  A verified 2014-format refutation ends with that
validated empty-clause addition \cite[Sec.~3 and Fig.~1]{WetzlerHeuleHunt2014AuthorDRAT}.
All local assumption trails are reset between tests; they are not added to
the permanent database.
\index{DRAT!addition and deletion trace}\index{DRAT!empty-clause check}
\end{hbdefinition}

\begin{hbmethod}{def:p12-drat-forward-checker}{A complete forward checking procedure}
Parse the CNF and the finite ordered proof according to the declared format
version; retain literal signs, pivot order, and the deletion interpretation.
Set the active database to the parsed CNF.  Process each record in order.  A
deletion performs the declared removal or no-op.  An addition runs the RUP
test; on failure it runs every opposing-clause RAT test if a pivot exists.
Commit the addition only after acceptance.  Reject malformed records and
failed additions.  Accept UNSAT only on a verified empty-clause addition.
For a finite trace containing fresh variables, take \(X\) to include all
variables of the input and that trace; existence of a valuation of the input
is unchanged by arbitrary extension to this finite superset.

This deliberately explicit checking method is needed to expose the trusted
logical condition.  It is not a claim that the 2014 optimized implementation
literally performs this forward pass.  The source describes backward
dependency checking, conflict-based marking, core-first propagation, and
scanning active clauses for a failed-RUP RAT test
\cite[Sec.~5]{WetzlerHeuleHunt2014AuthorDRAT}.  Such optimizations need their
own implementation invariants; the theorem below proves the stated method,
not omitted code paths.
\end{hbmethod}

\begin{hbtheorem}{thm:p12-drat-trace-sound}{proved=annex:p12-drat-trace-sound}{A verified DRAT refutation certifies UNSAT}
The finite checking procedure terminates.  For every accepted prefix its
active database satisfies
\[
 \operatorname{Sat}(F)\ \Rightarrow\ \operatorname{Sat}(F^i_P).
\]
Consequently, if the procedure accepts a verified empty-clause addition,
the original finite CNF \(F\) is unsatisfiable.  The guarantee also holds
for a multiset database with single-copy deletion, and for declared deletion
no-ops, because those operations remove constraints or leave them unchanged.
It does not require trust in the solver's search or in the correctness of
unaccepted log records.
\end{hbtheorem}

\begin{hbexample}{def:p12-drat-native-refutation}{The complete 2014 certificate example}
The source's Figure~1 input has variables \(a,b,c,d\), represented by the
integers \(1,2,3,4\), and these ten clauses:
\[
\begin{array}{lll}
 D_1=\{a,b,\neg c\},&D_2=\{\neg a,\neg b,c\},&D_3=\{\neg a,\neg b,\neg c\},\\
 D_4=\{b,c,\neg d\},&D_5=\{\neg b,\neg c,d\},&D_6=\{\neg a,\neg c,\neg d\},\\
 D_7=\{a,c,d\},&D_8=\{\neg a,b,d\},&D_9=\{a,\neg b,\neg d\},\\
 D_{10}=\{\neg a,b,\neg d\}.&&
\end{array}
\]
Its complete four-record proof is
\begin{verbatim}
-1 0
d -1 2 4 0
2 0
0
\end{verbatim}
For the first addition \(\{\neg a\}\), assuming \(a\) leaves no unit
clause, so RUP fails.  RAT has exactly three opposing clauses, \(D_1,D_7,D_9\).
Their pivot-retaining tests and complete conflict witnesses are:
\begin{itemize}
 \item \(\{\neg a,b,\neg c\}\): assume \(a,\neg b,c\).
       Clause \(D_6\) forces \(\neg d\), which falsifies \(D_8\).
 \item \(\{\neg a,c,d\}\): assume \(a,\neg c,\neg d\).
       Clause \(D_8\) forces \(b\), which falsifies \(D_2\).
 \item \(\{\neg a,\neg b,\neg d\}\): assume \(a,b,d\).
       Clause \(D_6\) forces \(\neg c\), which falsifies \(D_2\).
\end{itemize}
The first record therefore passes RAT.  The second deletes the present clause
\(D_8\).  To check the third, assume \(\neg b\) in the remaining
database, which contains the accepted unit \(\neg a\).  Then \(D_1\)
forces \(\neg c\), \(D_7\) forces \(d\), and \(D_4\) is false.
Thus the unit \(b\) passes RUP.  For the final empty-clause test no temporary
assumptions are added.  The retained units \(\neg a,b\) make \(D_9\)
force \(\neg d\); \(D_5\) then forces \(\neg c\), and \(D_7\)
is false.  The last record is verified, so the original ten-clause input is
UNSAT \cite[Fig.~1]{WetzlerHeuleHunt2014AuthorDRAT}.
The extra clauses \(D_3,D_{10}\) of this exact 2014 input have not been
silently replaced by the eight-clause input used in the 2016 note.
\end{hbexample}

\begin{hbexample}{def:p12-drat-failed-refutation}{A failed log with a countermodel}
For the input \(F=\{\{\neg p\}\}\), encode \(p\) by the variable
index \(1\) and consider the proposed records
\begin{verbatim}
1 0
0
\end{verbatim}
The input is SAT, witnessed by \(p=0\).  The attempted addition \(\{p\}\)
is not RUP: its sole assumption \(\neg p\) satisfies the input.  Nor is it
RAT: its opposing clause is \(\{\neg p\}\), giving the test clause
\(\{p\}\) again, which still fails RUP.  The checker must reject the
first record.  If an incorrect checker inserted it anyway, the pair of units
\(p,\neg p\) would make the final empty-clause test succeed; that conflict
would certify only the corrupted database, not the original input.
This independently worked failure case identifies the missing universal RAT
obligation and demonstrates why checking only the final conflict is unsound.
\end{hbexample}

\section{Construction, completeness, and logging obligations}

\begin{hbdefinition}{def:p12-drat-truth-tree}{An explicit finite certificate construction}
List the input variables as \(x_1,\ldots,x_n\).  For a prefix assignment
\(s\) to \(x_1,\ldots,x_k\), let its blocking clause \(B_s\) contain
the literal opposite to its value at each assigned variable.  Thus \(B_s\)
is false exactly on extensions of \(s\), and \(B_{\varnothing}\) is the
empty clause.  If \(F\) is unsatisfiable, add all full-assignment blocking
clauses, then the shorter prefix blocking clauses in decreasing prefix length,
ending with \(B_{\varnothing}\).  Each internal prefix has the two child
clauses \(B_s\cup\{x_{k+1}\}\) and
\(B_s\cup\{\neg x_{k+1}\}\) already present.  No deletions or fresh
variables are required.  This is a finite mathematical construction of a log,
not a performance prescription for production solvers.
\end{hbdefinition}

\begin{hbtheorem}{thm:p12-drat-drup-complete}{proved=annex:p12-drat-drup-complete}{Finite-CNF DRUP certificate completeness}
Every unsatisfiable finite CNF has a finite verified refutation using only RUP
additions.  The construction in Definition~\ref{def:p12-drat-truth-tree}
uses at most \(2^{n+1}-1\) additions and ends with a verified empty clause,
including the case \(n=0\).  Hence the mathematical DRUP subsystem and the
stated DRAT checking core are refutationally complete.  This existential
construction gives no polynomial bound on certificate size and no termination
guarantee for an arbitrary solver's heuristic search.
\end{hbtheorem}

A solver emitting a refutation must log every database change on which a later
certificate step depends, preserve the pivot as the first literal of each
non-RUP addition, and synchronize the declared deletion semantics with the
checker.  The solver's in-memory reason clauses or successful search are not
substitutes for the log.  Conversely a SAT claim is checked by a valuation of
the original input, not by a successful RAT prefix.  A timeout, interrupted
log, rejected proof, or log with no verified empty clause establishes no UNSAT
claim.  Completeness of the proof system, successful proof production, soundness
of the mathematical checker, and correctness of a particular checker binary
are four different obligations.

\section{Dated format variants and the trusted boundary}
\label{sec:p12-drat-format-versions}

\begin{hbdefinition}{def:p12-drat-format-2016}{The separately dated 2016 format conventions}
Heule's version~1 format note, arXiv:1610.06229v1, has these additional
conventions \cite[Syntax of Input Files; Clause Syntax Restrictions; Clause
Deletion Details]{Heule2016DRATFormatV1}.  Literals are nonzero signed
integers with positive variable indices at most \(2^{31}-1\); clauses end
at \texttt{0}, and a deletion has prefix \texttt{d}.
The input header gives an upper bound on the input variable indices and the
exact number of input clauses; every occurring index must be within that bound
and the clause count must match.  Both input and proof clauses exclude
duplicate literals and complementary pairs.  Here \(F\) denotes the input
as a \emph{multiset}, rather than the 2014 clause set:
two identical clause records are two copies, and a deletion removes only one
copy.  An absent-clause deletion is ignored with a warning.  The mathematical
trace semantics removes one present copy; the described DRAT-trim
implementation additionally ignores deletions of unit clauses and warns.
Acceptance is on a validated empty addition; records after it are ignored.
These are distinct parser and operational policies, not a new proof of RAT.

Removing a copy or doing nothing preserves satisfiability in the direction
needed by Theorem~\ref{thm:p12-drat-trace-sound}.  But a checker that ignores
unit deletions is checking its own modified active database; it need not accept
the same traces as literal single-copy deletion.  For example from
\(\{\{p\}\}\), deleting \(\{p\}\) makes adding \(\{\neg p\}\)
a vacuous RAT step; retaining the unit makes that same addition fail.
RUP is monotone under retention of extra clauses, whereas RAT also introduces
new opposing-clause obligations.  Therefore deletion policy is an observable
format boundary even when each policy is sound for its own accepted traces.
\end{hbdefinition}

\subsection{The separately dated binary serialization}
\label{sec:p12-drat-binary-2016}

The 2016 note also supplies a binary serialization: an addition begins with
the byte for \texttt{a}, a deletion with the byte for \texttt{d}, and a zero
byte ends the clause.  In this serialization paragraph \(l\) is a nonzero
signed integer literal code, not the abstract literal of the earlier semantic
definitions.  This code maps to \(2l\) for
\(l>0\) and \(-2l+1\) for \(l<0\); the resulting unsigned integer is
written in least-significant-first seven-bit groups, with a continuation bit
on all but the last byte \cite[Sec.~1]{Heule2016DRATFormatV1}.  This changes
representation, not the logical redundancy criterion.  A claim about a binary
checker additionally requires correct decoding, range handling, record
termination, and reconstruction of the same pivot and clause.

\begin{proofboundary}{def:p12-drat-trusted-boundary}{What independent checking establishes}
The mathematical theorem establishes UNSAT of the exact parsed CNF whenever
the specified procedure accepts its certificate.  Applying it to an executable
checker requires a correct parser and correspondence to that procedure, plus
the declared execution path: executable/compiler/runtime or a verified
extraction and its stated assumptions.  The independent checker need not trust
the SAT solver, its heuristics, or a proof-producing preprocessor.  It still
needs the exact input and certificate bytes.  A proof of the CNF says nothing
about the intended informal problem until the encoding's implication or
equisatisfiability theorem is supplied, as required by
Method~\ref{def:p12-encoding-audit}.  Independence of a checker is not itself
a theorem of checker correctness; the 2014 paper's timings and coverage of
then-known solving techniques are not universal complexity bounds.
\end{proofboundary}

The later LRAT interface adds clause identifiers and propagation/RAT hints
that a checker verifies; it does not permit trusting a hint merely because an
elaborator emitted it \cite[Secs.~1, 3]{CruzFilipeEtAl2016LRATV1}.  The
separately identified September~2021 FRAT preprint uses optional annotations,
clause identifiers, and final-database information for communication from a
solver to an elaborator \cite[Secs.~1--2]{BaekCarneiroHeule2021FRATV1}.
Neither format nor its verified-checker results is attributed to the original
2014 DRAT-trim paper.  A DRAT-to-LRAT or FRAT-to-LRAT pipeline can leave its
translator untrusted only when the final checker checks the produced proof
against the exact original input under its own proved specification.

RAT also does not preserve ordinary or projected model counts.  With a declared
single variable \(p\), the empty database has two models; adding the vacuous
RAT unit \(\{p\}\) leaves one.  Projecting onto \(\{p\}\) likewise
changes the two projected valuations to one.  Thus counting, projected
counting, enumeration, optimization, and QBF certificates need their own
semantic preservation conditions and checking formats.  They are not supplied
by the existence-of-a-model invariant proved here, nor by relabelling an
ordinary DRAT log.
